\section{2026 年 11 月核验增补}

本节按论文首次提交时间排列。每项方法附有可核验的外部入选依据。

% METHOD: 2402.03723
\paragraph{Rig3DGS：从手持单目视频构建可控制的肖像}
2024.02~\cite{2402_03723}。Rig3DGS 将头像高斯保存在规范空间，用可学习形变将头姿与表情映射到各帧。FLAME 网格的邻近顶点提供形变先验，学习到的修正让规范高斯沿表情运动，同时保留头发和眼镜等细节的表示。

论文在手机短视频的留出帧上分别评价固定相机时改变表情和头姿，以及固定头姿时改变视角与表情。基线可控制的自由度不同，原文没有在每个设置使用全部方法。作者指出强非均匀照明下的建模和主体明显移动仍有限。作者单位包含 Adobe Research。

\textbf{入选依据} 论文署名单位包括 Adobe Research（\href{https://arxiv.org/abs/2402.03723}{论文原文}）。

% METHOD: 2403.11589
\paragraph{UV Gaussians：在同一 UV 空间学习人体网格形变与高斯纹理}
2024.03~\cite{2403_11589}。UV Gaussians 从已标定同步的多视图视频建立可驱动人体模型。Mesh U-Net 根据姿态预测模板网格位移，Gaussian U-Net 在相同 UV 空间预测颜色、不透明度和形状等高斯属性，再依据细化网格将高斯带动到三维空间。网格提供随姿态变化的几何约束，高斯纹理保留衣褶等外观细节。

论文分别评估新视角和新姿态。补充协议使用 64 个训练视角、8 个测试视角、300 个训练帧和 100 个新姿态测试帧；为配合不同对照还分别拟合了 SMPL-X 和 SMPL。这样的采集规模和模板条件不能外推到无标定单目视频。腾讯 AI Lab 的署名是入选依据。

\textbf{入选依据} 论文署名单位包括 Tencent（\href{https://arxiv.org/html/2403.11589}{论文原文}）。

% METHOD: 2607.23861
\paragraph{动态显式发丝头像：分别建模头发与上半身运动}
2026.07~\cite{2607_23861}。该工作用多相机视频分别重建头部/上身与具有显式纤维结构的头发。发丝高斯受局部角速度、加速度和相对重力历史条件控制；与直接预测整个头部高斯形变相比，头发有自己的可解释组织和运动支路。

作者与多视图头像以及物理仿真基线比较发丝外观和运动。发丝与面部采用不同表示，在运动、分割或跟踪不准确时，接合处仍可能出现渲染伪影。作者 Marc Pollefeys 同时列有 ETH Zürich 与 Microsoft 单位。

\textbf{入选依据} 论文署名单位包括 Microsoft（\href{https://arxiv.org/html/2607.23861v1}{论文原文}）。

% METHOD: 2607.29633
\paragraph{OASIS：在单张手部图像上显式处理自遮挡}
2026.07~\cite{2607_29633}。OASIS 用手部几何把可见图像证据映射到手部表示，再以可见性条件化的点像交互为遮挡区域传播上下文。网格表面的特征提供局部非刚性关节组织，高斯用于生成可动画的单图手部头像。

论文在 InterHand2.6M 和随手输入比较，并分别展示直接预测和针对新手细化的设置。姿势、相机或掩码严重失准仍会损害纹理，遮住的表面通常依赖先验。作者仓库提供训练、评估和已经公开的权重入口。

\textbf{入选依据} \href{https://github.com/ivyyy77/oasis}{公开实现}已找到与论文对应的代码文件。

% METHOD: 2608.01726
\paragraph{G-Skin：以生成姿态图像为高斯骨架绑定提供监督}
2026.08~\cite{2608_01726}。G-Skin 接收已有静态高斯和给定骨架，将二维生成模型产生的多姿态图像用作蒙皮监督，在测试对象上优化高斯随骨骼运动的权重。几何正则约束相邻位置的变形，减弱生成图像不准确时的绑定漂移。

作者用 50 个类人和 35 个非人形的生成资产评价。关联部位已经粘连或道具与身体拓扑融合时，局部连续性可能促成错误绑定。此流程仍需逐对象优化，不能称为一步预测骨架或蒙皮权重。

\textbf{入选依据} 论文署名单位包括 Huawei（\href{https://arxiv.org/abs/2608.01726}{论文原文}）。

% METHOD: 2608.20335
\paragraph{4DAnyone：用骨架条件与跨视图合成支持单目人体四维重建}
2026.08~\cite{2608_20335}。4DAnyone 从单目视频恢复人体骨架，用骨架作为稀疏几何条件生成多视角视频。参考上下文压缩减少视角增加时的外观条件开销，目标视图的分组在去噪中动态关联，以减小各组之间的结构漂移。生成结果随后用于高斯人体重建。

测试覆盖 DNA-Rendering 和 DyMVHumans，并检查多视角合成与下游四维高斯结果。论文的失败例显示，远离骨架的宽松衣物可能跨视图不一致，错误骨架也可能被一致地传播到全部新视角。

\textbf{入选依据} 论文署名单位包括 Ant Group（\href{https://arxiv.org/abs/2608.20335}{论文原文}）。

% METHOD: 2608.23410
\paragraph{以逐尺度 Transformer 生成标准人脸视角再进行高斯提升}
2026.08~\cite{2608_23410}。该方法用逐尺度 Transformer 从一幅或多幅无已知相机位姿的条件图生成六个指定视角，再用现有像素对齐高斯提升模型重建三维头像。新方法主要位于生成标准视角的前级，三维提升承担跨视角重建与渲染。

论文在 PSGS 和 Ava-256 分析生成与重建，PSGS 的 ArcFace 指标并非所有对比中占优。训练集中性表情且无配饰，对新表情、帽子和背景颜色溢出仍有限制。原文列有多名 Meta 署名作者，首作者说明工作在 Meta Reality Labs 实习期间进行。

\textbf{入选依据} 论文署名单位包括 Meta Reality Labs（\href{https://arxiv.org/html/2608.23410v1}{论文原文}）。

% METHOD: 2610.26898
\paragraph{EGGHead：在固定原语预算内重建单图头部高斯}
2026.10~\cite{2610_26898}。EGGHead 学习按头部语义区域分配有限高斯，并借助 Voronoi 划分推断稀疏原语的局部形状。轻量对称模块补全遮挡区域，推断无需为每个新对象再执行逐个优化。

论文在不重合的 FFHQ 和 CelebA-HQ 图像上评价。模型未将纹理拆成反照率与照明，不能用于重打光；不均匀光线会影响对称补全，完全正脸时后脑仍可能不完整。

\textbf{入选依据} 论文署名单位包括 Qualcomm（\href{https://arxiv.org/abs/2610.26898}{论文原文}）。
